\documentclass[11pt]{article}
\usepackage[margin=1in]{geometry}
\usepackage{graphicx}
\usepackage{booktabs}
\usepackage{amsmath}
\usepackage{float}
\usepackage{caption}
\usepackage[round,authoryear]{natbib}
\usepackage[hidelinks]{hyperref}
\captionsetup{font=small}
\graphicspath{{../_output/}}
\title{Credit-Market Sentiment and the Business Cycle:\\ A Replication of \citet{lopezsalido2017credit}}
\author{Fernando Raffo Pedraza \and Bangjie Xu}
\date{\today}

\begin{document}
\maketitle

\begin{abstract}
We replicate the core results of \citet{lopezsalido2017credit}, who show that elevated
credit-market sentiment forecasts a slowdown in real activity about two years ahead. We
describe the data, present summary statistics of the underlying series, and reproduce every
table and figure across the replication sample, an extended sample, and an Aaa-spread variant.
\end{abstract}

\section{Overview}

\citet{lopezsalido2017credit} argue that credit-market sentiment---an unusually narrow
Baa--Treasury spread combined with an elevated share of low-quality (high-yield) bond
issuance---is a leading indicator of the business cycle. When credit is ``frothy'' in this
sense, spreads and issuer quality subsequently mean-revert and real activity slows roughly two
years later; wide spreads and depressed low-quality issuance tend to precede a recovery.

We reproduce the paper's core results from our own reconstruction of the underlying data, then
push on them in three directions: extending the sample through 2025, substituting the
Aaa--Treasury spread for the Baa--Treasury spread, and applying the fitted sentiment model out
of sample to the 2020--2022 COVID cycle. Section~\ref{sec:data} covers the data and the
judgement calls its reconstruction required, and Section~\ref{sec:successes} closes with where
we matched the published results and where we did not.

\section{Data Sources}
\label{sec:data}

We combine three public data sources, re-vintaged and spliced where necessary, to reconstruct
every series used in the replication:

\begin{itemize}
\item \textbf{FRED} \citep{fred}: Baa/Aaa yields, Treasury yields, GDP. Several series are a
newer vintage and must be stitched together across multiple FRED series. Where the appendix is
silent about which underlying series to use, we made documented judgement calls.
\item \textbf{Shiller} \citep{shiller2015data}: long-run market data and the P/E10 ratio.
\item \textbf{Greenwood--Hanson high-yield share} \citep{greenwood2013issuer}, pulled from the
HBS workbook \citep{hbsics} and spliced with Mergent FISD \citep{mergentfisd}, accessed via
WRDS; a newer vintage than the paper.
\end{itemize}

\section{Summary Statistics}
\label{sec:summary-stats}

Four processed series feed every regression in this report: the FRED credit-spread and
macroeconomic series, the Greenwood--Hanson high-yield issuance share, and the Shiller
equity-valuation series. We establish the scale and behavior of each one before it appears in
a regression table.

Over the FRED monthly sample, both the Baa- and Aaa--Treasury spreads sit well inside their
extremes for most of the period, with the dispersion driven by a handful of crisis episodes.

\begin{table}[H]
\centering
\caption{\label{tab:summary-credit-spreads}Summary statistics: Monthly FRED data (credit
spreads, yields, and recession share).}
\resizebox{\textwidth}{!}{%
\input{../_output/summary_statistics_credit_spreads.tex}
}
\end{table}

Both spreads widen sharply around NBER recessions (shaded), which is what makes them useful
cyclical indicators throughout the report.

\begin{figure}[H]
\centering
\includegraphics[width=0.85\textwidth]{summary_statistics_credit_spreads.pdf}
\caption{\label{fig:summary-credit-spreads}Credit spreads (Baa/Aaa--Treasury).}
\end{figure}

The annual FRED series cover real GDP, population, per-capita GDP, inflation, and long-run
corporate and Treasury yields---the levels from which the growth series and annual-frequency
regressors are built.

\begin{table}[H]
\centering
\caption{\label{tab:summary-gdp-growth}Summary statistics: Annual FRED data (GDP,
population, inflation, and yields/spreads).}
\resizebox{\textwidth}{!}{%
\input{../_output/summary_statistics_gdp_growth.tex}
}
\end{table}

Year-over-year growth in real GDP per capita is volatile, with the deepest contractions around
the Great Depression, World War II, and the 2008 and 2020 recessions. This is the outcome
variable every regression table forecasts.

\begin{figure}[H]
\centering
\includegraphics[width=0.85\textwidth]{summary_statistics_gdp_growth.pdf}
\caption{\label{fig:summary-gdp-growth}GDP growth.}
\end{figure}

The high-yield issuance share is the ingredient that separates \citet{lopezsalido2017credit}
from a pure credit-spread model, and its wide range reflects how much low-quality issuance
swings across credit cycles.

\begin{table}[H]
\centering
\caption{\label{tab:summary-hy-share}Summary statistics: High-yield issuance share.}
\resizebox{\textwidth}{!}{%
\input{../_output/summary_statistics_hy_share.tex}
}
\end{table}

The color break in Figure~\ref{fig:summary-hy-share} marks where the Greenwood--Hanson series
is spliced onto the Mergent FISD reconstruction used for the post-2008 extension and the case
study.

\begin{figure}[H]
\centering
\includegraphics[width=0.85\textwidth]{summary_statistics_hy_share.pdf}
\caption{\label{fig:summary-hy-share}High-yield issuance share.}
\end{figure}

Shiller's data supply the S\&P Composite price and dividend levels, the cyclically-adjusted
P/E ratio, and total log returns---the equity-valuation series that competes against the
credit spread in Table~\ref{tab:table1-replication}.

\begin{table}[H]
\centering
\caption{\label{tab:summary-cape}Summary statistics: Cyclically-adjusted P/E (CAPE) and
forward equity returns.}
\resizebox{\textwidth}{!}{%
\input{../_output/summary_statistics_cape.tex}
}
\end{table}

Higher CAPE valuations are followed by weaker ten-year S\&P returns (correlation shown on the
chart). That relationship is the standard motivation for treating equity valuation as a
serious rival to the credit spread as a predictor.

\begin{figure}[H]
\centering
\includegraphics[width=0.65\textwidth]{summary_statistics_cape.pdf}
\caption{\label{fig:summary-cape}Cyclically-adjusted P/E (CAPE) and forward equity returns.}
\end{figure}

\section{Replication Sample (1929--2015)}
\label{sec:replication}

We start on the paper's original 1929--2015 sample, working through the long-run behavior of
the credit spread, the horse race between the spread and equity returns, the
sentiment-augmented regression that splits the spread into mean-reverting and ``froth''
components, and the resulting link between sentiment and subsequent growth.

The Baa--Treasury spread widens sharply into every recession, which is what makes its lagged
changes a natural candidate predictor of future growth.

\begin{figure}[H]
\centering
\includegraphics[width=0.85\textwidth]{figure_1_replication.pdf}
\caption{\label{fig:spread-replication}Baa--Treasury credit spread with NBER recessions
shaded, 1929--2015 (paper's Figure I).}
\end{figure}

Table~\ref{tab:table1-replication} runs the paper's central horse race: one-year-ahead GDP
growth on the lagged change in the credit spread, against the same specification using lagged
equity returns. The credit-based specification's markedly higher $R^2$ is the paper's evidence
that credit prices carry predictive information equity prices do not.

\begin{table}[H]
\centering
\caption{\label{tab:table1-replication}One-year-ahead GDP growth on the lagged
credit-spread change vs.\ equity returns, 1929--2015 (paper's Table I).}
\input{../_output/table_1_replication.tex}
\end{table}

Table~\ref{tab:table2-replication} splits the spread into a mean-reverting component and a
froth component predicted by the lagged high-yield share, then uses both to forecast future
spread changes. Their joint significance is why the paper treats issuer quality, not just the
spread level, as carrying independent predictive content.

\begin{table}[H]
\centering
\caption{\label{tab:table2-replication}Sentiment-based predictive regressions with the
auxiliary decomposition, 1929--2015 (paper's Table II).}
\input{../_output/table_2_replication.tex}
\end{table}

Figure~\ref{fig:sentiment-replication} is the paper's central result in one picture: the
fitted sentiment measure slopes downward against subsequent GDP growth, so elevated sentiment
precedes weaker growth.

\begin{figure}[H]
\centering
\includegraphics[width=0.65\textwidth]{figure_2_replication.pdf}
\caption{\label{fig:sentiment-replication}Credit-market sentiment vs.\ subsequent growth,
1929--2015 (paper's Figure II).}
\end{figure}

\section{Extended Sample (1929--2025)}
\label{sec:extended}

Re-estimating everything on data through 2025 adds a decade of history, including the 2020
pandemic recession, to relationships originally fit on 1929--2015.

The spread continues to widen around each subsequent recession, 2020 included, which also
confirms the series was constructed consistently across the extension.

\begin{figure}[H]
\centering
\includegraphics[width=0.85\textwidth]{figure_1_extended.pdf}
\caption{\label{fig:spread-extended}Baa--Treasury credit spread with NBER recessions shaded,
extended to 2025.}
\end{figure}

The extended horse race asks whether the credit spread's edge over equity returns belongs to
the 1929--2015 window specifically or is a more durable property of the data.

\begin{table}[H]
\centering
\caption{\label{tab:table1-extended}Table I regressors re-estimated on the extended sample,
1929--2025.}
\input{../_output/table_1_extended.tex}
\end{table}

The extended decomposition adds another decade---including the post-2008 high-yield market
captured by the FISD splice---to test whether the froth and mean-reversion components remain
jointly significant.

\begin{table}[H]
\centering
\caption{\label{tab:table2-extended}Table II regressors re-estimated on the extended sample,
1929--2025.}
\input{../_output/table_2_extended.tex}
\end{table}

The sentiment--growth relationship survives the additional years without noticeable dilution.

\begin{figure}[H]
\centering
\includegraphics[width=0.65\textwidth]{figure_2_extended.pdf}
\caption{\label{fig:sentiment-extended}Credit-market sentiment vs.\ subsequent growth,
extended to 2025.}
\end{figure}

\section{Extension: Aaa Credit Spread (Replication Sample, 1929--2015)}
\label{sec:aaa-replication}

Swapping the Aaa--Treasury spread for the Baa--Treasury spread tests how much of the paper's
result depends on using medium-grade rather than high-grade credit.

The higher-grade spread still widens around recessions, but visibly less than Baa: Aaa issuers
are far less exposed to credit-quality concerns in a downturn.

\begin{figure}[H]
\centering
\includegraphics[width=0.85\textwidth]{figure_1_aaa_replication.pdf}
\caption{\label{fig:spread-aaa-replication}Aaa--Treasury credit spread with NBER recessions
shaded, 1929--2015.}
\end{figure}

Re-running the horse race on Aaa shows whether the credit-spread advantage over equity returns
is specific to Baa-grade credit.

\begin{table}[H]
\centering
\caption{\label{tab:table1-aaa-replication}Table I regressors using the Aaa--Treasury
spread, 1929--2015.}
\input{../_output/table_1_aaa_replication.tex}
\end{table}

The same question applies to the sentiment decomposition: does the froth/mean-reversion split
still carry predictive content for a spread with less credit-quality exposure?

\begin{table}[H]
\centering
\caption{\label{tab:table2-aaa-replication}Table II regressors using the Aaa--Treasury
spread, 1929--2015.}
\input{../_output/table_2_aaa_replication.tex}
\end{table}

The sentiment--growth slope is flatter than its Baa counterpart, consistent with Aaa sentiment
carrying weaker predictive content.

\begin{figure}[H]
\centering
\includegraphics[width=0.65\textwidth]{figure_2_aaa_replication.pdf}
\caption{\label{fig:sentiment-aaa-replication}Credit-market sentiment built on the Aaa
spread vs.\ subsequent growth, 1929--2015.}
\end{figure}

\section{Extension: Aaa Credit Spread (Extended Sample, 1929--2025)}
\label{sec:aaa-extended}

Combining both robustness checks re-estimates the Aaa specification on the extended
1929--2025 sample.

The smoother Aaa cyclical pattern persists into recent data.

\begin{figure}[H]
\centering
\includegraphics[width=0.85\textwidth]{figure_1_aaa_extended.pdf}
\caption{\label{fig:spread-aaa-extended}Aaa--Treasury credit spread with NBER recessions
shaded, extended to 2025.}
\end{figure}

The extended Aaa horse race confirms the same-signed result in
Table~\ref{tab:table1-aaa-replication} is not an artifact of the shorter window.

\begin{table}[H]
\centering
\caption{\label{tab:table1-aaa-extended}Table I regressors using the Aaa--Treasury spread,
extended sample.}
\input{../_output/table_1_aaa_extended.tex}
\end{table}

The extended Aaa decomposition stays consistent with the replication-sample estimates in
Table~\ref{tab:table2-aaa-replication}.

\begin{table}[H]
\centering
\caption{\label{tab:table2-aaa-extended}Table II regressors using the Aaa--Treasury spread,
extended sample.}
\input{../_output/table_2_aaa_extended.tex}
\end{table}

Through 2025 the Aaa sentiment--growth relation weakens further, but it does not reverse sign.

\begin{figure}[H]
\centering
\includegraphics[width=0.65\textwidth]{figure_2_aaa_extended.pdf}
\caption{\label{fig:sentiment-aaa-extended}Credit-market sentiment built on the Aaa spread
vs.\ subsequent growth, extended to 2025.}
\end{figure}

\section{Case Study: The Sentiment Signal Out of Sample, 2020--2022}
\label{sec:case-study}

Reconstructing the post-2008 high-yield issuance share from Mergent FISD \citep{mergentfisd}
lets us take the Table II first stage---estimated on 1929--2015 and held fixed---to the COVID
cycle. The test is genuinely out of sample: the FISD data did not exist when the original
estimates were fit.

Figure~\ref{fig:case-study} compares the predicted change in the Baa--Treasury spread against
the realized change over 2020--2022, using the full two-ingredient sentiment signal and,
separately, the spread-reversion leg alone. With the first stage pinned at its 1929--2015
estimates, agreement between predicted and realized changes is evidence the relationship is
not an artifact of the training sample.

\begin{figure}[H]
\centering
\includegraphics[width=0.85\textwidth]{case_study_covid_oos.pdf}
\caption{\label{fig:case-study}Out-of-sample application of the sentiment first stage to
2020--2022: predicted versus realized change in the credit spread.}
\end{figure}

\section{Successes and Challenges}
\label{sec:successes}

The long-run credit spread reproduces cleanly, and the paper's central credit-beats-equity
result holds on both the replication and extended samples. The sentiment decomposition in
Table~\ref{tab:table2-replication} and the resulting downward sentiment--growth relationship
in Figure~\ref{fig:sentiment-replication} also come through as in the paper, and both survive
extending the sample to 2025 and substituting the Aaa--Treasury spread---though the Aaa-based
results are visibly weaker, consistent with the paper's own framing of Baa as the more
informative grade.

The main difficulties were on the data side: sourcing and splicing the Greenwood--Hanson
high-yield share across its pre- and post-2009 vintages, and matching the paper's
\citet{newey1987simple} standard errors exactly. Series-vintage differences (FRED revisions,
the Greenwood--Hanson/FISD splice) and our own judgement calls where the paper's appendix is
silent explain the small remaining gaps between our numbers and the published ones. The
specific calls are described in Section~\ref{sec:data}.

\bibliographystyle{jpe}
\bibliography{references}

\end{document}
